% ============================================================
\chapter{Points Fixes dans les Ensembles Ordonn\'es}
\label{ch:ordonnes}
% ============================================================

En 1928, Bronislaw Knaster et Alfred Tarski d\'emontrent un r\'esultat d'une simplicit\'e et d'une port\'ee extraordinaires~: toute application croissante sur un treillis complet poss\`ede un point fixe. Mieux encore, l'ensemble de \emph{tous} les points fixes forme lui-m\^eme un treillis complet. Ce th\'eor\`eme n'utilise ni m\'etrique, ni topologie, ni continuit\'e~: seule la structure d'ordre intervient. Il trouvera des applications spectaculaires en informatique th\'eorique, o\`u il fonde la s\'emantique d\'enotationnelle des programmes r\'ecursifs (les travaux de Dana Scott dans les ann\'ees 1970), et en logique, o\`u il sous-tend la th\'eorie des d\'efinitions inductives. Avec les it\'er\'es transfinis de Bourbaki--Witt et le th\'eor\`eme de Kleene pour les $\omega$-cpo, ce chapitre explore la troisi\`eme grande famille de th\'eor\`emes de points fixes --- celle fond\'ee sur l'ordre.

\begin{intuition}
Les th\'eor\`emes de points fixes dans les ensembles ordonn\'es constituent la troisi\`eme
grande famille de r\'esultats, apr\`es les th\'eor\`emes m\'etriques (Banach) et
topologiques (Brouwer, Schauder). Ici, aucune structure m\'etrique ou topologique n'est
requise~: seule la structure d'ordre intervient. Le r\'esultat fondamental est le
th\'eor\`eme de Knaster--Tarski, qui affirme que toute application croissante sur un
treillis complet poss\`ede un point fixe, et que l'ensemble des points fixes forme
lui-m\^eme un treillis complet.
\end{intuition}

% ------------------------------------------------------------
\section{Rappels sur les ensembles ordonn\'es}
% ------------------------------------------------------------

\begin{definition}[Ordre partiel et treillis]
\index{ordre partiel}
Un \emph{ensemble partiellement ordonn\'e} (poset) est un couple $(P, \leq)$ o\`u $\leq$
est une relation r\'eflexive, antisym\'etrique et transitive. Un \emph{treillis} est un
poset dans lequel toute paire $\{a, b\}$ admet un supremum $a \vee b$ et un infimum
$a \wedge b$. Un treillis est \emph{complet} si toute partie admet un supremum et un
infimum.
\end{definition}

\begin{exemple}[Treillis complets fondamentaux]
\begin{itemize}
  \item $(\mathcal{P}(X), \subseteq)$ pour tout ensemble $X$.
  \item $([0,1], \leq)$ avec l'ordre usuel.
  \item Le treillis des sous-espaces vectoriels d'un espace vectoriel.
  \item Le treillis des relations d'\'equivalence sur un ensemble.
\end{itemize}
\end{exemple}

\begin{definition}[Application croissante et d\'ecroissante]
\index{application croissante}
Soit $(P, \leq)$ un poset. Une application $f\colon P \to P$ est \emph{croissante}
(ou monotone) si $x \leq y \Rightarrow f(x) \leq f(y)$. Elle est \emph{inflationnaire}
si $x \leq f(x)$ pour tout $x$, et \emph{d\'eflationnaire} si $f(x) \leq x$ pour tout $x$.
\end{definition}

% ------------------------------------------------------------
\section{Th\'eor\`eme de Knaster--Tarski}
% ------------------------------------------------------------

\begin{theoreme}[Knaster--Tarski]
\label{thm:knaster-tarski}
\index{th\'eor\`eme de Knaster--Tarski}
Soit $(L, \leq)$ un treillis complet et $f\colon L \to L$ une application croissante.
Alors~:
\begin{enumerate}
  \item $f$ admet un plus petit point fixe~:
    $\mathrm{lfp}(f) = \bigwedge \{x \in L : f(x) \leq x\}$.
  \item $f$ admet un plus grand point fixe~:
    $\mathrm{gfp}(f) = \bigvee \{x \in L : x \leq f(x)\}$.
  \item L'ensemble $\mathrm{Fix}(f) = \{x \in L : f(x) = x\}$ est un treillis complet
    (pour l'ordre induit).
\end{enumerate}
\end{theoreme}

\begin{proof}
Posons $\mathrm{Pre}(f) = \{x \in L : f(x) \leq x\}$ (les pr\'e-points fixes) et
$a = \bigwedge \mathrm{Pre}(f)$.

\emph{\'Etape 1~:} $f(a) \leq a$. En effet, pour tout $x \in \mathrm{Pre}(f)$, on a
$a \leq x$, donc $f(a) \leq f(x) \leq x$. Ainsi $f(a)$ minore $\mathrm{Pre}(f)$,
d'o\`u $f(a) \leq a$.

\emph{\'Etape 2~:} $a \leq f(a)$. De l'\'etape~1, $f(a) \leq a$, donc par croissance,
$f(f(a)) \leq f(a)$. Ainsi $f(a) \in \mathrm{Pre}(f)$, d'o\`u $a \leq f(a)$.

\emph{\'Etape 3~:} $f(a) = a$ par les \'etapes 1 et 2.

\emph{Plus petit point fixe~:} si $b$ est un point fixe, alors $b \in \mathrm{Pre}(f)$,
d'o\`u $a \leq b$.

\emph{Treillis complet des points fixes~:} Soit $S \subseteq \mathrm{Fix}(f)$. D\'efinissons
$L_S = \{x \in L : x \geq s \; \forall s \in S \text{ et } f(x) \leq x\}$. Alors
$L_S$ est non vide ($\top \in L_S$) et $\bigwedge L_S$ est un point fixe de $f$ qui est
le supremum de $S$ dans $\mathrm{Fix}(f)$. Dualement, $\mathrm{Fix}(f)$ admet les infima.
\end{proof}

\begin{exemple}[Application en logique]
Soit $X$ un ensemble et $\Phi\colon \mathcal{P}(X) \to \mathcal{P}(X)$ un op\'erateur
croissant. Le plus petit point fixe de $\Phi$ est le plus petit ensemble $A$ tel que
$\Phi(A) = A$. Si $\Phi$ est un op\'erateur de cons\'equence logique, $\mathrm{lfp}(\Phi)$
est l'ensemble des formules d\'eductibles.
\end{exemple}

% ------------------------------------------------------------
\section{Th\'eor\`eme de Bourbaki--Witt}
% ------------------------------------------------------------

\begin{theoreme}[Bourbaki--Witt]
\label{thm:bourbaki-witt}
\index{th\'eor\`eme de Bourbaki--Witt}
Soit $(P, \leq)$ un ensemble ordonn\'e dans lequel toute cha\^ine admet un majorant, et
soit $f\colon P \to P$ une application inflationnaire ($x \leq f(x)$ pour tout $x$). Si
$P$ poss\`ede un plus petit \'el\'ement $\bot$, alors $f$ admet un point fixe.
\end{theoreme}

\begin{proof}
Consid\'erons l'ensemble $A$ de toutes les cha\^ines $C$ de $P$ telles que~:
\begin{itemize}
  \item $\bot \in C$~;
  \item $C$ est stable par $f$~;
  \item si $C' \subseteq C$ est une cha\^ine, tout majorant de $C'$ dans $P$ qui est dans
    $C$ est bien un majorant.
\end{itemize}
On d\'efinit $C_0 = \bigcap_{C \in A} C$. On montre que $C_0$ est une cha\^ine bien
ordonn\'ee. Soit $m$ un majorant de $C_0$ dans $P$ (qui existe par hypoth\`ese). Alors
$f(m) \geq m$. Si $f(m) = m$, c'est un point fixe. Sinon, $C_0 \cup \{m\}$ serait une
cha\^ine strictement plus grande, ce qui contredit la minimalit\'e de $C_0$.
\end{proof}

\begin{remarque}
Le th\'eor\`eme de Bourbaki--Witt est \'equivalent au lemme de Zorn dans ZF + axiome du
choix. Sans l'axiome du choix, l'\'enonc\'e reste vrai pour les fonctions inflationnaires
explicitement d\'efinies.
\end{remarque}

% ------------------------------------------------------------
\section{Application \`a l'analyse formelle de concepts}
% ------------------------------------------------------------

\begin{definition}[Contexte formel]
\index{analyse formelle de concepts}
Un \emph{contexte formel} est un triplet $(G, M, I)$ o\`u $G$ est un ensemble d'objets,
$M$ un ensemble d'attributs, et $I \subseteq G \times M$ une relation d'incidence ($gIm$
signifie que l'objet $g$ poss\`ede l'attribut $m$).
\end{definition}

\begin{definition}[Op\'erateurs de d\'erivation]
Pour $A \subseteq G$ et $B \subseteq M$, on d\'efinit~:
\[
  A' = \{m \in M : \forall g \in A, \; gIm\}, \qquad
  B' = \{g \in G : \forall m \in B, \; gIm\}.
\]
\end{definition}

\begin{proposition}[Connexion de Galois]
Les op\'erateurs $A \mapsto A'$ et $B \mapsto B'$ forment une \emph{connexion de Galois}
entre $(\mathcal{P}(G), \subseteq)$ et $(\mathcal{P}(M), \supseteq)$~:
\begin{enumerate}
  \item $A_1 \subseteq A_2 \Rightarrow A_2' \subseteq A_1'$~;
  \item $A \subseteq A''$~;
  \item $A' = A'''$.
\end{enumerate}
\end{proposition}

\begin{definition}[Concept formel]
Un \emph{concept formel} est un couple $(A, B)$ avec $A \subseteq G$, $B \subseteq M$,
$A' = B$ et $B' = A$. L'ensemble $A$ est l'\emph{extension} et $B$ l'\emph{intension}
du concept.
\end{definition}

\begin{theoreme}[Treillis des concepts]
\label{thm:treillis-concepts}
L'ensemble des concepts formels d'un contexte $(G, M, I)$, ordonn\'e par
$(A_1, B_1) \leq (A_2, B_2) \Leftrightarrow A_1 \subseteq A_2$ (ou de mani\`ere
\'equivalente $B_2 \subseteq B_1$), est un treillis complet.
\end{theoreme}

\begin{proof}
L'op\'erateur de fermeture $\varphi\colon \mathcal{P}(G) \to \mathcal{P}(G)$ d\'efini par
$\varphi(A) = A''$ est croissant, extensif ($A \subseteq A''$) et idempotent
($A'' = A''''$). Les ensembles ferm\'es (points fixes de $\varphi$) forment un treillis
complet, et correspondent bijectivement aux concepts formels.
\end{proof}

\begin{exemple}
Consid\'erons le contexte avec $G = \{1, 2, 3\}$, $M = \{a, b, c\}$ et
$I = \{(1,a), (1,b), (2,b), (2,c), (3,a), (3,c)\}$. L'extension $\{1,2\}' = \{b\}$,
et le concept correspondant est $(\{1,2\}, \{b\})$.
\end{exemple}

% ------------------------------------------------------------
\section{Points fixes et op\'erateurs de fermeture}
% ------------------------------------------------------------

\begin{definition}[Op\'erateur de fermeture]
\index{op\'erateur de fermeture}
Un \emph{op\'erateur de fermeture} sur un poset $(P, \leq)$ est une application croissante,
extensive et idempotente $\varphi\colon P \to P$.
\end{definition}

\begin{theoreme}
Soit $(L, \leq)$ un treillis complet et $\varphi\colon L \to L$ un op\'erateur de fermeture.
Alors $\mathrm{Fix}(\varphi)$ est un treillis complet (avec les suprema de $L$ referm\'es
par $\varphi$, et les infima inchang\'es).
\end{theoreme}

\begin{proof}
Les points fixes de $\varphi$ sont les \'el\'ements ferm\'es, c'est-\`a-dire les $x$ tels
que $\varphi(x) = x$. Pour $S \subseteq \mathrm{Fix}(\varphi)$, l'infimum dans
$\mathrm{Fix}(\varphi)$ est $\bigwedge_L S$ (car $\varphi(\bigwedge_L S) \leq
\varphi(s) = s$ pour tout $s \in S$, d'o\`u $\varphi(\bigwedge_L S) \leq \bigwedge_L S$,
et par extensivit\'e l'\'egalit\'e). Le supremum dans $\mathrm{Fix}(\varphi)$ est
$\varphi(\bigvee_L S)$.
\end{proof}

\begin{formules}[title=\textbf{Points fixes dans les ensembles ordonn\'es}]
\begin{center}
\renewcommand{\arraystretch}{1.3}
\begin{tabular}{|l|l|l|}
\hline
\textbf{Th\'eor\`eme} & \textbf{Hypoth\`eses} & \textbf{Conclusion} \\
\hline
Knaster--Tarski & Treillis complet, $f$ croissante & $\mathrm{Fix}(f)$ treillis complet \\
Bourbaki--Witt & Cha\^ines born\'ees, $f$ inflat. & Point fixe \\
Op. fermeture & Treillis complet, $\varphi$ fermeture & $\mathrm{Fix}(\varphi)$ treillis complet \\
\hline
\end{tabular}
\end{center}
\end{formules}

% ------------------------------------------------------------
\section{Exercices}
% ------------------------------------------------------------

\begin{exercice}
Montrer que le th\'eor\`eme de Knaster--Tarski implique le th\'eor\`eme de
Cantor--Bernstein~: si $f\colon A \hookrightarrow B$ et $g\colon B \hookrightarrow A$ sont
des injections, alors il existe une bijection $h\colon A \to B$.
\emph{Indication~: appliquer le th\'eor\`eme \`a $\varphi(S) = A \setminus g(B \setminus f(S))$
sur $(\mathcal{P}(A), \subseteq)$.}
\end{exercice}

\begin{exercice}
Soit $(L, \leq)$ un treillis complet et $f\colon L \to L$ croissante. Montrer que si
$a \leq f(a)$ et $f(b) \leq b$ avec $a \leq b$, alors $f$ admet un point fixe dans
l'intervalle $[a, b]$.
\end{exercice}

\begin{exercice}
Donner un exemple d'ensemble ordonn\'e (non complet) et d'application croissante sans
point fixe.
\end{exercice}

\begin{exercice}
Soit $G = (\{1,2,3,4\}, E)$ un graphe et $f\colon \mathcal{P}(\{1,2,3,4\}) \to
\mathcal{P}(\{1,2,3,4\})$ d\'efini par $f(S) = S \cup \{v : \exists u \in S, \; (u,v) \in E\}$.
Calculer $\mathrm{lfp}(f)$ et $\mathrm{gfp}(f)$ pour le graphe
$E = \{(1,2), (2,3), (3,4)\}$ avec $f$ d\'emarrant de $\{1\}$.
\end{exercice}

\begin{exercice}[Connexion de Galois et points fixes]
\index{connexion de Galois}
Soit $(f, g)$ une connexion de Galois entre deux treillis complets $L_1$ et $L_2$
($f \dashv g$). Montrer que $g \circ f$ est un op\'erateur de fermeture sur $L_1$ et
$f \circ g$ un op\'erateur de noyau sur $L_2$.
\end{exercice}

\begin{exercice}
Construire le treillis des concepts du contexte formel avec $G = \{\text{chien},
\text{chat}, \text{poisson}\}$, $M = \{\text{poils}, \text{nage}, \text{4 pattes}\}$ et
la relation d'incidence naturelle.
\end{exercice}
